\chapter{Untouchability and Caste: Past, Present and Future}

\section{Diwali, Civil Religion and the Little/Great Tradition}

\begin{ExampleBox}
\begin{itemize}
\item \textbf{Why do we celebrate Diwali?} Different religions celebrate it for different reasons: the Hindu Brahmanical view (the return of Rama with Sita and Lakshmana); the \textbf{Jain} (Mahavira attained moksha); the \textbf{Sikh} (Guru Har Gobind Singh released from Jahangir's Gwalior Fort --- \textbf{Bandi Chhod Divas}); and the Bengali (Kali/Durga puja, not Rama).
\item At the end of the day, everyone has a different private religion, but Diwali becomes a \textbf{public religion for everyone} --- a \textbf{civic religion} (Robert Bellah's idea): good over evil, light over darkness, knowledge over ignorance.
\item The \textbf{diyas are totems} --- symbols whose interpretation depends on one's \textbf{social location} (mechanical solidarity). The many different reasons are the \textbf{little tradition}, while the unified 'Diwali' we think of is the \textbf{great tradition}.
\end{itemize}
\end{ExampleBox}

\section{Michael Moffatt: The Replicability Principle}

\begin{ExampleBox}
\begin{itemize}
\item \textbf{Michael Moffatt} conducted an \textbf{ethnographic study} (about \textbf{14 months}) of two untouchable castes --- the \textbf{Paraiyars} and the \textbf{Pallars} --- in the village of \textbf{Andavur}, \textbf{Chingleput} district, Tamil Nadu (about 50 miles from Madras).
\item Social composition: the \textbf{untouchables are 32\%} of the village; the \textbf{dominant caste is the Reddiar} (a peasant caste, over 50\%). The village is divided into \textbf{Ur} (the elite area) and \textbf{Cheri} (the inferior area) --- settlement patterns (pur, ur, nadu, puram).
\end{itemize}
\end{ExampleBox}

\begin{KeyConcept}
\begin{itemize}
\item The \textbf{book view} of untouchability is a \textbf{'disjunctive model'} --- treating touchables and untouchables as two separate, polar, binary categories (a structuralist perspective). Moffatt says it is time to change this.
\item His key finding: \textbf{untouchables are both excluded and included} --- on \textbf{exclusion} they \textbf{replicate}; on \textbf{inclusion} they \textbf{subordinate}. This is the \textbf{replicability (replicatory) principle}.
\item Because they are excluded, for their own survival they \textbf{replicate the same hierarchy and structure} of the mainstream caste society within themselves.
\end{itemize}
\end{KeyConcept}

\section{Replication Within Untouchables}

\begin{KeyConcept}
\begin{itemize}
\item Moffatt found \textbf{three strategies of replication}:
\item \textbf{1. Replication of the caste hierarchy} --- within the untouchables there is a hierarchy just like the mainstream:
\item \textbf{Valluvar Pandarams} --- the '\textbf{Brahmins of the Harijans}', the \textbf{purest among the impure} (they conduct life-cycle rituals, marriage, death and even astrology).
\item \textbf{Paraiyars} --- \textbf{numerically preponderant} among the untouchables (about 80\%) --- a \textbf{replica of the dominant Reddiar}.
\item \textbf{Vannan} --- washermen, specifically washing the clothes of girls at first menstruation, and arranging funeral pyres, and cutting hair.
\item \textbf{Chakkiliyar} --- leather workers (working with dead cattle).
\item \textbf{Kurivikaran} --- bird-catchers who \textbf{eat crows}, placed lowest among the untouchables (the most impure).
\item \textbf{2. Caste divided into sub-castes, each sub-caste endogamous} --- the Paraiyars themselves are divided into sub-castes, and endogamy is maintained within each.
\item \textbf{3. Replication of Hindu religion} --- excluded from the Ur's religious functions, the untouchables form \textbf{their own gods}: the Ur's local deity is \textbf{Seliyaman}, while the Cheri's is its replica \textbf{Mariamman}.
\item This is the same strategy A. M. Shah saw in the \textbf{Garo} of Gujarat --- adding Brahminical surnames (Joshi, Vyas) and conducting Brahmin rituals.
\end{itemize}
\end{KeyConcept}

\begin{KeyConcept}
\textbf{Subordination (on inclusion)}: when allowed to take part in an upper-caste ritual, they \textbf{accept their low status}. The Uber/Ola analogy (Michael Burawoy's \textbf{'work games'}): a platform worker works for two companies at once as a \textbf{defence mechanism} --- work turns into a game that liberates from exploitation.
\end{KeyConcept}

\section{Subordination and the Critique of the Disjunctive Model}

\begin{KeyConcept}
\begin{itemize}
\item Moffatt's theory is a \textbf{critique of the disjunctive models of untouchability} (which stressed the contrast between untouchables and the higher castes).
\item When purity/impurity is found \textbf{even amongst the impure}, it shows \textbf{ideological homogeneity} --- Dumont's purity/impurity is the ideology of the whole Hindu civilization, found in every sphere.
\item \textbf{Dumont}: the ideology of Hindu civilization is \textbf{highly non-competitive} --- the impure accept their position (the 'tire never competes with the engine'). Untouchables accept their inequality so completely that they \textbf{replicate purity/impurity within themselves}. Hindu inequality is a \textbf{special form of inequality not found anywhere else}.
\end{itemize}
\end{KeyConcept}

\section{Caste: Past, Present and Future}

\begin{KeyConcept}
\begin{itemize}
\item \textbf{'Is caste dead / weakened / strengthened?'} --- caste is \textbf{not a quantitative, measurable variable}; it cannot be said to have 'strengthened' or 'weakened' (that is TV-journalist talk). The change is \textbf{qualitative --- a change in form}.
\item \textbf{Ghurye}: in traditional society caste was highly flexible; as India evolved it became rigid; with the British it became \textbf{both flexible and rigid}.
\item \textbf{Srinivas}: earlier \textbf{vertical solidarity}, now \textbf{horizontal solidarity}; hierarchy earlier based on \textbf{ritual observance}, now on \textbf{non-ritual (political, economic, numerical) criteria}; the \textbf{new avatar of caste = the OBCs} (the 20th-century avatar, based on social/economic backwardness, not purity/impurity).
\item \textbf{Rajni Kothari}: earlier \textbf{casteization of politics}, now \textbf{politicization of caste}.
\item \textbf{D. L. Sheth}: \textbf{structural differentiation} within caste --- it had ritual functions, now it also has non-ritual political functions --- the \textbf{secularization of caste}.
\item \textbf{Beteille}: earlier the \textbf{old caste}; now the \textbf{new caste replaces the old caste} (class did not replace caste).
\end{itemize}
\end{KeyConcept}

\section{Caste Outside India}

\begin{KeyConcept}
\begin{itemize}
\item \textbf{What caste is to Indians, race is to outsiders.} The racial classificatory system has \textbf{Caucasoids as superior} and \textbf{Negroids as inferior} --- just as the (Risley's) racial theory of caste ranked the fair at the top and the dark at the bottom.
\item \textbf{Oliver Cox}: \textbf{caste is just like race} --- race is '\textbf{coloured caste}' (endogamy within race, as within caste; each race has its own lifestyle).
\item \textbf{Ghurye}: the \textbf{class structure of Britain is the exact replica of the caste structure of India} --- we have casteism, they have classism.
\item Sociologists look not at the \textbf{manifest body} but the \textbf{latent soul} --- caste operates in different forms/avatars outside India too.
\end{itemize}
\end{KeyConcept}

\section{From Caste Public to Bio-Public}

\begin{KeyConcept}
\begin{itemize}
\item In the \textbf{contemporary (post-pandemic)} context, untouchability has evolved: before the pandemic there was \textbf{vertical untouchability} (the impure below the pure, even within untouchables). After the pandemic there is \textbf{horizontal untouchability} --- social distancing is now maintained between the \textbf{COVID-positive and the non-positive}, irrespective of caste.
\item This is the move from a \textbf{'caste public' to a 'bio-public'} (concern has shifted from \textbf{ritual hygiene to medical hygiene}) --- a concept given by the sociologist \textbf{Badri Narayan}.
\item Social distancing existed before and exists now --- only the \textbf{form has changed} (again, not strengthen/weaken, but a change in form).
\end{itemize}
\end{KeyConcept}

\begin{QuickRevision}
\begin{itemize}
\item Moffatt (Paraiyars \& Pallars, Andavur, Chingleput): untouchables both excluded and included --- on exclusion they replicate, on inclusion they subordinate (replicability principle).
\item Replication: (1) hierarchy within untouchables (Valluvar Pandarams = 'Brahmins of Harijans'; Paraiyars = dominant-replica; Vannan, Chakkiliyar, Kurivikaran at the bottom); (2) sub-castes, each endogamous; (3) own gods (Seliyaman $\rightarrow$ Mariamman); same as Garo (A. M. Shah).
\item Subordination = accepting low status on inclusion (Burawoy's 'work games').
\item Critique of disjunctive models; ideological homogeneity; Dumont's 'non-competitive' ideology.
\item Caste past/present/future = qualitative change, not quantitative: Ghurye (flexible$\rightarrow$rigid$\rightarrow$both), Srinivas (vertical$\rightarrow$horizontal; ritual$\rightarrow$secular; OBC = new avatar), Kothari (casteization$\rightarrow$politicization), D. L. Sheth (secularization), Beteille (new caste replaces old).
\item Caste outside India: race (Caucasoid/Negroid); Oliver Cox (caste = race; 'coloured caste'); Ghurye (British class = Indian caste).
\item Post-pandemic: vertical $\rightarrow$ horizontal untouchability; caste public $\rightarrow$ bio-public (Badri Narayan); ritual hygiene $\rightarrow$ medical hygiene.
\end{itemize}
\end{QuickRevision}

